\documentclass[../../../include/open-logic-section]{subfiles}

\begin{document}

\olfileid{sth}{z}{power}
\olsection{घातसंच}

आता आणखी एक स्वयंसिद्धक पाहू:

\begin{axiom}[घातसंच]
कोणत्याही संच $A$ साठी $\Pow{A} = \Setabs{x}{x \subseteq A}$ हा संच अस्तित्वात असतो.
\[
	\forall A \exists P \forall x(x \in P \liff (\forall z \in x)z \in A)
\]
\end{axiom}

याचे समर्थन अगदी सरळ आहे. $A$ हा संच टप्पा $S$ वर तयार झाला असे मानू. मग \stagesacc{} नुसार $A$ चे सर्व सदस्य $S$ च्या आधी उपलब्ध होते. म्हणून विभक्तीकरणाच्या समर्थनाप्रमाणे विचार केल्यास $A$ चा प्रत्येक उपसंच टप्पा $S$ पर्यंत तयार झालेला असतो. ते सर्व उपलब्ध असल्याने $S$ नंतरच्या कोणत्याही टप्प्यावर त्यांचा एक संच तयार करता येतो. \stagessucc{} नुसार $S$ हा शेवटचा टप्पा नाही, म्हणून असा पुढील टप्पा आहेच. त्यामुळे $\Pow{A}$ अस्तित्वात असतो.

घातसंचाच्या स्वयंसिद्धकाचा एक चांगला परिणाम असा:

\begin{prop}\label{thm:Products}
कोणतेही संच $A, B$ दिल्यास त्यांचा कार्तीय गुणाकार $A \times B$ अस्तित्वात असतो.
\end{prop}

\begin{proof}
घातसंचाचे स्वयंसिद्धक आणि \olref[sth][z][pairs]{prop:pairsconsequences} यांनुसार $\Pow{\Pow{A \cup B}}$ अस्तित्वात असतो. म्हणून विभक्तीकरणानुसार पुढील संच अस्तित्वात असतो:
\[
	C = \Setabs{z \in \Pow{\Pow{A \cup B}}}{(\exists x \in A)(\exists y \in B) z = \tuple{x, y}}.
\]
आता कोणतेही $x \in A$ आणि $y \in B$ दिल्यास \olref[sth][z][pairs]{prop:pairsconsequences} नुसार $\tuple{x, y}$ अस्तित्वात असतो. शिवाय $x, y
\in A \cup B$ असल्याने $\{x\}, \{x, y\} \in \Pow{A \cup B}$ आणि $\tuple{x,y} \in \Pow{\Pow{A \cup B}}$ मिळते. म्हणून $A \times B = C$.
\end{proof}

या सिद्धतेत घातसंच आणि विभक्तीकरण एकत्र काम करतात. त्यात आश्चर्य नाही. विभक्तीकरण नसते, तर घातसंचाचे तत्त्व फारसे \emph{सामर्थ्यवान} ठरले नसते. कारण कोणते उपसंच अस्तित्वात असतात हे विभक्तीकरण सांगते; त्यामुळे प्रत्येक घातसंच किती ``भरलेला'' असेल, हेही त्यावर अवलंबून असते.

\begin{prob}
कोणतेही संच $A, B$ दिल्यास दाखवा: (i) प्रांत $A$ आणि व्याप्ती $B$ असलेल्या सर्व संबंधांचा संच अस्तित्वात असतो; आणि (ii) $A$ पासून $B$ कडे जाणाऱ्या सर्व फलनांचा संच अस्तित्वात असतो.
\end{prob}

\begin{prob}
$A$ हा संच आणि $\sim$ हा $A$ वरील तुल्यता संबंध मानू. $\sim$ नुसार $A$ वरील तुल्यतावर्गांचा संच, म्हणजे $\equivclass{A}{\sim}$, अस्तित्वात असतो हे सिद्ध करा.
\end{prob}

\end{document}
